% GUID v6, appendices.
\section{Reproducibility}\label{app:repro}

Every number of this paper is printed by one of the scripts in Table~\ref{tab:scripts}. Each script prints OK or FAIL for each of its
checks and exits with code 0 only if all pass. The one exception is the data-confrontation script, which exits with code 1 on purpose
because three scored confrontations fail at 95\% (Section~\ref{sec:data}). \texttt{python3 docs/guid-v6/reproduce.py} runs all of them and
prints the table again from the scripts' own output.

All scripts are deterministic. They fix their random seeds, pin the linear-algebra libraries to one thread, and read no clock except to
report their own runtime. Two runs of each gave identical output files apart from that runtime field. The runtimes in the table are for
one core of a shared machine and are not results. The scripts need Python~3 with numpy and scipy; \texttt{checks.py} also needs sympy,
\texttt{holographic\_code.py} needs stim~1.16, and the scripts that draw figures need matplotlib. None of them needs a network connection.

\begin{center}\small
\begin{longtable}{@{}P{0.26\linewidth}P{0.17\linewidth}P{0.15\linewidth}P{0.09\linewidth}P{0.22\linewidth}@{}}
\caption{The scripts behind the paper. The first two are \texttt{docs/guid-v6/checks.py} and \texttt{docs/qg-verdict/assembly.py}; the
others are in \texttt{docs/simulation/}.}\label{tab:scripts}\\
\toprule
Script & Sections & Checks & Runtime & Labels \\
\midrule
\endfirsthead
\toprule
Script & Sections & Checks & Runtime & Labels \\
\midrule
\endhead
\bottomrule
\endfoot
\texttt{checks.py} & Part~I & 125 of 125 & seconds & \chk{Q1}--\chk{Q20} \\
\texttt{assembly.py} & \ref{sec:holo} & 18 of 18 & seconds & \chk{A1}--\chk{A5} \\
\texttt{hosted\_benchmarks.py} & \ref{sec:hosted} & 53 of 53 & 8 min & \chk{hosted:B1}--\chk{hosted:B4} \\
\texttt{holographic\_code.py} & \ref{sec:holocode} & 77 of 77 & 8 min & \chk{holo:H0}--\chk{holo:H6} \\
\texttt{gravity\_simulated.py} & \ref{sec:gravsim}, \ref{sec:bmv} & 110 of 110 & 3 min & \chk{gravsim:G1}--\chk{gravsim:G4} \\
\texttt{bmv\_witness.py} & \ref{sec:bmv} & 34 of 34 & 1 min & \chk{witness:B1}--\chk{witness:B4} \\
\texttt{sparc\_rotation.py} & \ref{sec:sparc} & 39 of 39 & 1 min & \chk{sparc:S0}--\chk{sparc:S3} \\
\texttt{remnant\_history.py} & \ref{sec:remnant} & 39 of 39 & 2.5 min & \chk{remnant:R1}--\chk{remnant:R6} \\
\texttt{predictions\_vs\_data.py} & \ref{sec:data} & 6 of 6; 28 of 31 scored & 4 s & \chk{data:D1.1}--\chk{data:D9.6} \\
\texttt{standard\_model.py} & \ref{sec:sm} & 32 of 32 & 5 s & \chk{sm:SM1}--\chk{sm:SM5} \\
\texttt{exotic\_hadrons.py} & \ref{sec:exotics} & 80 of 80 & 1 min & \chk{exotics:A}--\chk{exotics:H} \\
\texttt{tally.py} & \ref{sec:tally} & 6 of 6 & 1 s & \chk{tally:total} \\
\end{longtable}
\end{center}

\paragraph{Data.} The rotation-curve analyses read the SPARC tables of Lelli, McGaugh and Schombert~\cite{lelli2016}. Their licence is not
stated, so they are not redistributed with this paper. Both scripts check the files by SHA-256 before use and stop if a digest differs:
\begin{itemize}[nosep]
\item \texttt{MassModels\_Lelli2016c.mrt}:\\ {\footnotesize\texttt{9108994b12cc401b94a1768beca61c53ec354779385c9c9cc571049f3043244c}};
\item \texttt{SPARC\_Lelli2016c.mrt}:\\ {\footnotesize\texttt{5aa0501f6b0d881fa579030e315e7b5b6ef561a5bd3a07472f9929c7e5728243}}.
\end{itemize}
Pass their folder with \texttt{--data <folder>} or the environment variable \texttt{SPARC\_DATA\_DIR}. All other data enter the scripts
as published values, each with its source in the script and in Sections~\ref{sec:data} to~\ref{sec:exotics}.

\paragraph{Results files.} Each simulation script writes its numbers to a JSON file and a readable account to a Markdown file of the
same name, both in \texttt{docs/simulation/results/}; its figures go to \texttt{docs/simulation/figures/}. The JSON files are the reference for the labels of the form \chk{script:label} in the
text; a label that names a whole analysis step (such as \chk{sparc:S0} or \chk{remnant:R6}) points to the step in the script and its
results note.

\section{The checks of Part I}\label{app:checks}

\texttt{python3 docs/guid-v6/checks.py} (numpy, scipy, sympy) prints 125 checks, all of which pass.
\begin{itemize}[nosep]
\item Q1: Planck units, the existence energy, Unruh, Hawking, Bekenstein--Hawking, Jacobson's entropy density, and the size of the quantum
term of~\eqref{eq:eft}.
\item Q2: the lapse.
\item Q3: no-cloning and broadcasting.
\item Q4: classical against coherent mediation, 300 random couplings each.
\item Q5: the Bose et al.\ geometry.
\item Q6: the MICROSCOPE bound on the information scalar, at the one-standard-deviation level and at 95\%, and the Higgs-share test.
\item Q7: Brans--Dicke and Cassini.
\item Q8: polarisations and Weyl components in $D$ dimensions.
\item Q9: the single-graviton threshold.
\item Q10: Kerr horizon, ergosphere, Christodoulou, efficiencies.
\item Q11: Hawking radiation of real holes.
\item Q12: the Hubble mass at its Schwarzschild threshold.
\item Q13: flatness, the de Sitter floor and the onset of acceleration.
\item Q14: the tensor spectrum, the BICEP/Keck bound and Starobinsky's $r$.
\item Q15: what every trapped remnant must satisfy, and why the known neutrinos do not (with the DESI~DR2 mass bound).
\item Q16: $\alpha$ as self-interaction strength, the electron's lengths, gravity against electromagnetism, and the running of $\alpha$.
\item Q17: the speed of gravity.
\item Q18: remnants trapped in the field: the stochastic remnant's rotation curve and acceleration scale, the contained remnant's Hawking
temperatures and evaporation, and the radiated remnant's energy density.
\item Q19: the two bridges, Bridge~A against the Tully--Fisher normalisation, Bridge~B's windows and Di\'osi--Penrose rates, and the
granular term at both smearing bounds.
\item Q20: the radiation-era clock and the time of the electroweak crossover, the W and Z masses from the Higgs vacuum, $\alpha$ at its
birth and today, and the proton lifetime for a unification at the start.
\end{itemize}
\texttt{python3 docs/qg-verdict/assembly.py} (numpy) prints the 18 checks of the holographic-code assembly. A1 to A5 in this paper are
its checks [1] to [5]; its check [6] is the special case of \chk{Q4} with CNOT and CZ couplings.
